\documentclass[a4paper,11pt]{article}

\usepackage{revy}
\usepackage[utf8]{inputenc}
\usepackage[T1]{fontenc}
\usepackage[danish]{babel}
\usepackage{amsmath,amssymb}

\revyname{MatRevy}
\revyyear{2026}
\version{2026-10-25}
\eta{3.5 minutter}

\title{Beskidt Matematik}
\author{Markus '20}

\begin{document}
\maketitle

\begin{roles}
\role{I}[Hugo] Instruktør
\role{S1}[Oliver] Skuespil
\role{S2}[Philip] lille person
\role{OE}[Markus] Onde Ernst
\role{N1}[Ninja] Ninja, stærk person
\role{N2}[Ninja] Ninja, stærk person
\end{roles}

\begin{sketch}
\scene{S1 og S2 er almindelige studerende. Onde Ernst er klædt som en tegneserie skurk ala Dracula. OE har magiske krafter og kan få matematiske ting til at flyve, som ninjaerne bærer. Evt. græskebogstaver eller udtryk}
\says{S1} Øj matintro er bare så spændende, men altså det med studieretningerne forvirrer mig.
\says{S2} Hvad forvirrer dig?
\says{S1}Jamen der er gymnasiel, statistik og så ren matematik?
\says{S2} Det giver da god mening?
\says{S1} Altså jaja, men nej.. altså ren matematik, hvorfor er det så rent? Findes der beskidt matematik???
\scene{Lys slukker, skarpt lysblink,*TORDEN SKRALD*}
\scene{Mørkt bag tæppe går til siden, OE træder frem…}
\says{OE} hmhmmh var der nogen der sagde… beskidt matematik?
\says{S1 og S2} øøhhh
\says{OE} var det dig!?\act{Peger på S1}
\says{S1}\act{Peger på S2}
\says{OE} \act{STORMER mod S2 stopper ham med kraften og tager fat i kraven og ved hjælp af to ninjaer bliver S2 løftet op i sin krave.}
\says{S2}\act{Fuld af rædsel i ansigtet.}
\says{S2}[Stammende] Hvad nej? Hvem er du.
\says{OE} hvem jeg er?? \act{slipper S2 som bliver hængende i luften}
\says{OE} JEG ER ONDE ERNST!
\scene{*TORDEN SKRALD*}
\says{OE}[Griber fat i S2 igen] men det egentlige spørgsmål er ikke hvad beskidt matematik er. Det er om i er værdige til at blive åbnet op til en ny verden af overtællelige dimensioner.
\says{OE} \act{styrer S2 ned (som the force)}
\scene{Ninjaer sætter S2 langsomt ned}
\says{S1 og S2} \act{ kigger på hinanden i frygt}
\says{S1-2} Det ved vi ikke om vi er
\says{OE} Nej selvfølgeligt ved I ikke det. Det er noget jeg vurderer.
\says{OE} S1 hvornår har du sidst gjort rent i snullet?
\says{S1} uhhh... aldrig?
\says{OE} hmm goodt… godt… S2! Kan du differentiere denne funktion? \act{Hiver en funktion frem?}
\says{S2} øhhh, men altså, den er jo ikke engang vel-defineret!
\says{OE} KORREKT!
\says{OE} Nu-Well, i viser lovende potientiale… Lad mig uddybe. Beskidt matematik, er alt det ufine. En stolt tradition udbygget af en lang række dygtige matematikkere tit overset af deres mere skole-rene tvillinger. Vi taler om Karl friedrich snavs, Wallodi Weibullshit og John von Nullermann. De har lavet det hårde arbejde, imens teoretikerne sidder i deres elfenbenstårn ude af trit med hvordan vi andre får verden til at dreje rundt.
\says{OE} \act{Står triumferende med armene til siden}
\says{S1} Vil det sige at Gauss og de andre har haft en ond tvillingebror?
\says{S2} Ja findes der også en Kolmo-go-røv?
\says{OE} Var din mund! Dette er ikke noget at spøge med. \act{Peger på S2}
\says{N1} {dækker S2’s mund}
\says{OE} \act{Går truende mod S2}
\says{S1}[Prøver at lede OE væk fra S2] Hey, nå men hvad lærer man så i beskidt matematik.
\says{N1} {stopper med at dække S2’s mund}
\says{OE}\act{grynter} I første uge vil i lære at dobbelt, tripple og quadruple indexere jeres summer,  i vil sørge for at fortegnene altid går op.

\says{S2} \act{Gisp}
\says{OE}[henkastende] Ja ja trivielt, men bare hør her
\scene{OE fører dem rundt i lokalet}
\says{OE} I vil skulle kunne formelen for partiel integration udenad og med lukkede øjne. 
\says{S1} Men det kan man da ikke?
\scene{OE fører dem videre}
\says{OE} Her er ingen kære mor; udvalgs aksiomet er ikke et valg! Det er tvang!
\says{S2} Men er det ikke farligt!?
\says{OE} Hahah… jo meget farligt, for de uværdige altså. I de rette hænder og med min vejledning kan i opnå uanede resultater! Ingen opgaver efterlades til LÆSEREN! Intet er TRIVIELT! Alt skal ses efter i sømmene og tjekkes til mindste detalje!
\scene{OE bliver mere og mere ophidset}
\scene{OE pruster og stønner af overanstrengelse}
\says{OE} Når i rene matematiker samvittighedsløst taster tal ind på sin lommeregner tror i bare svaret kommer af sig selv.
\scene{OE kigger bebrejdende ud på publikum}
\says{OE} NEEEJ det er mig der sidder derinde og regner efter! Slut jer til den beskidte side og før min arv videre. 
\says{S1} Det lyder da vildt kedeligt bare at sidde at tjekke indekser og lave regnestykker.
\says{OE} KEDELIGT?? Du forstår ikke engang betydningen af det ord. Dag ud og dag ind sidder jeg og tjekker ALLE induktionstrinnene - hver og et. Døgnet rundt. Når i bruger høje tal, en milliard, er der jo nogen der skal tælle dertil først! For ikke at nævne de overtællelige mængder! (Mere og mere desperat)
\scene{S1 og S2 flygter langsomt}
\says{OE} NEJ VENT! SLut jer til mig. Befri mig fra mine lænker.
\scene{Spot på Onde Ernst, der slået sætter sig ned.}
\says{OE} Aaahhhh. De forlader mig. Som alle andre. Årh. Jeg er så træt af at tælle.
\scene{Lys ned}
\end{sketch}
\end{document}
